% Generated draft: requires mathematical and source review.
\documentclass[11pt]{article}
\usepackage[margin=25mm]{geometry}
\usepackage{fontspec}
\setmainfont{Latin Modern Roman}
\usepackage{amsmath,amssymb,mathtools,unicode-math}
\setmathfont{Latin Modern Math}
\usepackage{enumitem,needspace,xurl,hyperref}
\setlength{\parindent}{0pt}
\setlength{\parskip}{5pt}
\setlength{\emergencystretch}{3em}
\newcommand{\N}{\mathbb{N}}
\newcommand{\Z}{\mathbb{Z}}
\newcommand{\Q}{\mathbb{Q}}
\newcommand{\R}{\mathbb{R}}
\newcommand{\C}{\mathbb{C}}
\newcommand{\F}{\mathbb{F}}
\newcommand{\A}{\mathbb{A}}
\newcommand{\PP}{\mathbb{P}}
\newcommand{\E}{\mathbb{E}}
\newcommand{\eps}{\varepsilon}
\newcommand{\dd}{\,\mathrm{d}}
\begin{document}
\section*{UnsolvedMath ID: 700242}
\textbf{Source problem code:} catalogue entry
\subsection{Definitions and mathematical statement}
Let \(v=v(x_1,\ldots,x_r)\) be a group word. For a group \(G\), define its set of \(v\)-values and its verbal subgroup by
\[
V_v(G)=\{v(g_1,\ldots,g_r):g_1,\ldots,g_r\in G\},
\qquad vG=\langle V_v(G)\rangle.
\]
Define \(\mathfrak K_v\) to be the class of groups \(G\) for which
\[
\exists n\in\Z_{>0}\ \forall h\in vG\ \exists (g_{ij})\in G^{n\times r}
\quad h=\prod_{i=1}^{n}v(g_{i1},\ldots,g_{ir}).
\]
Here \(n\) may depend on \(G\), with \(v\) fixed; it must work for every element of \(vG\). Since \(v(1,\ldots,1)=1\), allowing at most \(n\) factors gives the same condition as allowing exactly \(n\) factors. The supplied wording is followed literally: factors are word values, with inverse values admitted only when they are themselves word values. Whether a convention separately allowing inverses was intended requires source review.

For subgroups \(H\leq G\), write
\[
[H,H]=\langle a^{-1}b^{-1}ab:a,b\in H\rangle.
\]
A group \(G\) is soluble if its derived series, defined by \(G^{(0)}=G\) and \(G^{(j+1)}=[G^{(j)},G^{(j)}]\), satisfies \(G^{(d)}=\{1\}\) for some integer \(d\geq0\). It is finitely generated if it is generated by a finite subset.

\begin{enumerate}
\item[(a)] Characterize the group words \(v\) such that
\[
\text{for every finitely generated soluble group }G,
\qquad G\in\mathfrak K_v.
\]
No bound uniform over all such groups \(G\) is requested.
\item[(b)] Does the commutator word
\[
v(x,y)=[x,y]=x^{-1}y^{-1}xy
\]
satisfy the condition in part~(a)? Equivalently, is it true that for every finitely generated soluble group \(G\), there exists \(n\in\Z_{>0}\) such that every \(h\in G'=[G,G]\) has an expression
\[
h=\prod_{i=1}^{n}[a_i,b_i],\qquad a_i,b_i\in G?
\]
\end{enumerate}
\subsection{Short English statement}
Determine which group words have finite verbal width, in the supplied sense of products of word values, in every finitely generated soluble group. The bound may depend on the group. In particular, ask whether every finitely generated soluble group has finite commutator width.
\subsection{Sources}
\begin{itemize}
\item Local UnsolvedMath registry, supplied complete record with permanent ID \(700242\), exact source problem code \texttt{catalogue entry}, title \emph{Kourovka Notebook Problem 4.34}.
\item \emph{The Kourovka Notebook (Unsolved Problems in Group Theory)}, problem 4.34, attributed in the supplied record to M.\ I.\ Kargapolov, proposed in 1973 in the fourth issue.
\item The supplied background identifies the twenty-first-issue edition (2026), \texttt{arXiv:1401.0300v48}. The collection description names E.\ I.\ Khukhro and V.\ D.\ Mazurov as editors. These bibliographic details are supplied attribution and have not been independently checked.
\item The supplied registry labels the problem \emph{open}. This is a source status label, not an independently verified assessment of the literature.
\end{itemize}
\end{document}
